\section{Whole relative monomial stages under smooth submersions}

The original underlying set is arbitrary. In particular, none of the
following statements substitutes the whole ambient space for that set.
The ideal, complete ordered boundary and original marking are retained.

\begin{theorem}[Every original divisibility predicate and exponent along the original set]
\label{hir:smooth-relative-monomial-divisibility}
\uses{hir:smooth-relative-whole-cosupport,hir:submersion-local-section}
\lean{Hironaka.HironakaSmoothSubmersionRelativeMonomial.isDivNearAlong_preimage_iff,Hironaka.HironakaSmoothSubmersionRelativeMonomial.divExpAlong_preimage_eq}
For a genuine holomorphic submersion $f:Y^m\to X^n$, any original sets
$N,H\subset X$ and coherent original ideal $I$, every original nearby
divisibility predicate along $N$ is exactly equivalent to its counterpart
along $f^{-1}N$ and $f^{-1}H$. The actual original exponent is unchanged.
\end{theorem}
\begin{proof}
The actual holomorphic local section through each $y$ proves the exact
neighbourhood-filter identity $f_*(\mathcal N_y)=\mathcal N_{f(y)}$:
continuity gives one inclusion, and composition with the actual section
gives the reverse inclusion. The entire relative cosupport in every
natural degree is the actual preimage of its original counterpart.
Substitute this equality into the original nearby divisibility predicate,
and transport the eventual statement through the filter identity. The
underlying conditions are exactly membership in $N$ and $H$, so this
identifies every degree predicate, without regularity of either set.
The original exponent is defined by the first failed successor-degree
predicate, with its actual zero fallback when no failure exists. Equality
of the whole predicate family identifies both the first failure and the
fallback, proving the literal exponent equality.
\end{proof}
\leanok

\begin{theorem}[Entire indexed original relative exponent and invariant dictionary]
\label{hir:smooth-relative-monomial-whole-dictionary}
\uses{hir:smooth-relative-monomial-divisibility}
\lean{Hironaka.HironakaSmoothSubmersionRelativeMonomial.expDataAlong_preimage_eq,Hironaka.HironakaSmoothSubmersionRelativeMonomial.gammaAlong_preimage_dictionary,Hironaka.HironakaSmoothSubmersionRelativeMonomial.gammaAlong_same_dimension_preimage_eq}
The whole actual indexed exponent function is preserved for every original
ordered boundary entry. The original relative monomial invariant on the
preimage is computed by those same exponents and original through-indices
in the genuine source dimension. In equal ambient dimensions its whole
value is literally equal to the original value.
\end{theorem}
\begin{proof}
Preimage preserves the literal boundary-list length and every indexed
entry. For an index in range apply the preceding actual exponent equality
to precisely its original member. For an index outside the list, both
original functions have their actual value zero. Thus the entire indexed
functions agree, including historical empty and repeated entries.
Membership in each pulled member is literally membership of its image in
the same original entry; hence the whole through-index sets are equal.
Substitute both full equalities into the original definition of the
relative monomial invariant, retaining the genuine ambient dimension.
When the dimensions agree, this is exactly the original entire value,
including bottom and all colex ties.
\end{proof}
\leanok

\begin{theorem}[Whole original relative ordering across genuine ambient dimensions]
\label{hir:smooth-relative-monomial-ordering}
\uses{hir:smooth-relative-monomial-whole-dictionary,hir:submersion-monomial-invariant-order}
\lean{Hironaka.HironakaSmoothSubmersionRelativeMonomial.gammaAlong_preimage_le_iff,Hironaka.HironakaSmoothSubmersionRelativeMonomial.gammaAlong_preimage_eq_iff}
If the entire original ordered boundary is SNC, every comparison and
every equality between original relative monomial values at two actual
image points is preserved and reflected by the smooth preimage.
\end{theorem}
\begin{proof}
The actual submersion linear coordinates internally imply $n\le m$.
The entire original SNC coordinates internally bound the cardinality of
each genuine through-index set by $n$. The preceding whole indexed
dictionary supplies exactly the same feasible subsets in both dimensions.
The actual dimension-comparison theorem for the original monomial
invariant therefore identifies both weak comparisons, retaining empty
feasible families, bottom values and original colex ordering. Equality is
the conjunction of the two opposite comparisons. Neither dimension nor
cardinality bounds, invariant values or ordering certificates are inputs.
\end{proof}
\leanok

\begin{theorem}[Whole original relative marked maximum centre]
\label{hir:smooth-relative-monomial-whole-maximum}
\uses{hir:smooth-relative-monomial-ordering,hir:smooth-relative-whole-cosupport}
\lean{Hironaka.HironakaSmoothSubmersionRelativeMonomial.maxLocus_gammaAlong_markedPullback_eq}
For genuine surjective smooth $f$, a coherent original marked ideal and
entire original SNC boundary, the whole relative monomial maximum locus
on its actual relative cosupport pulls back exactly. The original
underlying set and marking are arbitrary.
\end{theorem}
\begin{proof}
The complete original relative cosupport dictionary gives exactly the
preimage set. A pulled point is maximal precisely when every other actual
pulled cosupport point has value at most its value. Each original
cosupport point has an actual preimage by genuine surjectivity; that
preimage lies in the pulled cosupport by the same whole dictionary.
Reflect every comparison to obtain maximality on the entire original
cosupport. Conversely original maximality preserves every comparison at
each actual pulled cosupport point. This proves equality of the entire
ambient maximum sets. It retains empty cosupports and arbitrary markings,
and uses no compactness, caller maximum or ambient-set certificate.
Higher invariant equality, canonical blowup base change and complete
canonical smooth functoriality remain open.
\end{proof}
\leanok
